\documentclass[10pt]{article}
\usepackage[margin=1in]{geometry}
\usepackage{amsmath,amssymb,amsthm}
\usepackage{hyperref}
\usepackage{graphicx} % for \resizebox
\usepackage{enumitem}
\setlist{itemsep=1pt,topsep=2pt}

% References (APA)
\usepackage[backend=biber,style=apa,natbib=true]{biblatex}
\addbibresource{references.bib}

\newtheorem{proposition}{Proposition}
\newtheorem{corollary}{Corollary}
\newtheorem{remark}{Remark}
\newcommand{\E}{\mathbb{E}}
\newcommand{\R}{\mathbb{R}}
\newcommand{\ip}[2]{\langle #1, #2 \rangle}

\title{\vspace{-2em}Bilinear GCIQL as a Low-Rank Factorization of Discounted Goal Reachability:\
Connections to VIP, Quasimetric RL, and Contrastive RL}
\author{}
\date{July 2026}

\begin{document}
\maketitle
\vspace{-3em}

\paragraph{Setup.}
Consider a goal-conditioned MDP with state space $\mathcal{S}$, deterministic dynamics (as in OGBench cube when the full state is observed), sparse reward $r(s;g) = \mathbf{1}[s{=}g] - 1$, termination upon reaching the goal, and discount factor $\gamma$. Let $d^*(s,g) \in {0,1,\dots,\infty}$ denote the minimum number of steps required to reach $g$ from $s$ under policies supported by the offline dataset. The optimal goal-conditioned value is
\begin{equation}
V^*(s,g) ;=; -\frac{1-\gamma^{d^*(s,g)}}{1-\gamma}, \qquad \gamma^{\infty} := 0 .
\label{eq:vstar}
\end{equation}
$d^*$ is a \emph{quasimetric}: $d^*(s,s)=0$ and $d^*(s,g) \le d^*(s,w) + d^*(w,g)$, while symmetry is not required (Wang et al.\ 2023, Thm.~1: quasimetrics are \emph{exactly} the class of optimal goal-reaching values). Bilinear GCIQL parametrizes $V(s,g) = \ip{\varphi(s)}{\psi(g)}/\sqrt{d}$ and $Q(s,a,g)=\ip{\varphi(s,a)}{\psi(g)}/\sqrt{d}$ and trains these functions using expectile temporal-difference learning. As the expectile parameter $\tau \to 1$, $V$ approaches $V^*$ over actions supported by the offline dataset. We make the statement that ``the critic learns a metric'' precise through three identities.

\section{The negated value is a quasimetric}

\begin{proposition}[Discounted quasimetric]\label{prop:qm}
Let $\phi_\gamma(x) := \frac{1-\gamma^{x}}{1-\gamma}$. Then $D_\gamma := -V^* = \phi_\gamma \circ d^*$ is a bounded quasimetric ($D_\gamma \le \tfrac{1}{1-\gamma}$).
\end{proposition}
\begin{proof}
$\phi_\gamma$ is increasing, concave, and satisfies $\phi_\gamma(0)=0$; therefore, it is subadditive, $\phi_\gamma(a+b)\le\phi_\gamma(a)+\phi_\gamma(b)$. Thus,
$\phi_\gamma(d^*(s,g)) \le \phi_\gamma(d^*(s,w)+d^*(w,g)) \le \phi_\gamma(d^*(s,w))+\phi_\gamma(d^*(w,g))$, and $D_\gamma(s,g)=0 \iff d^*(s,g)=0$. This is a metric-preserving transformation; the same argument appears implicitly in Wang et al.\ 2023.
\end{proof}

Therefore, the training target of bilinear GCIQL is already a quasimetric and does not require a logarithmic transformation. Equivalently, define the \textbf{discounted goal-reachability function}
\begin{equation}
K^*(s,g) := \gamma^{,d^*(s,g)} ;=; 1 + (1-\gamma),V^*(s,g) ;\in; [0,1].
\label{eq:kernel}
\end{equation}

\begin{corollary}[Composition inequality for discounted reachability]\label{cor:mult}
$K^*(s,s)=1$ and $K^*(s,g) ;\ge; K^*(s,w),K^*(w,g)$ for all $s,w,g$. This follows directly from the triangle inequality and $\gamma<1$.
\end{corollary}

\section{Bilinear GCIQL as a low-rank factorization of $K^*$}

Substituting the bilinear parametrization into \eqref{eq:kernel}, with $\tilde\varphi(s) := \big(1,; \tfrac{1-\gamma}{\sqrt d}\varphi(s)\big)$ and $\tilde\psi(g) := \big(1,; \psi(g)\big)$, gives
\begin{equation}
K(s,g) ;=; \ip{\tilde\varphi(s)}{\tilde\psi(g)}, \qquad \operatorname{rank} \le d+1 .
\label{eq:factorization}
\end{equation}
\emph{Bilinear GCIQL can represent the optimal discounted goal-reachability function exactly only when $\gamma^{d^*}$ admits a rank-$d$ factorization, and approximately when it admits an accurate rank-$d$ approximation.} This condition concerns the singular-value spectrum of the state-goal reachability matrix rather than the optimization procedure. The observed sharp performance transition on cube across representation dimensions (success $9.6\%$ at $d{=}5$ and $\sim 90\%$ at $d{=}50$) is consistent with the effective numerical rank of $K^*$ on state-goal pairs supported by the offline dataset. The discount factor controls how quickly reachability decreases with temporal distance: $K^* = e^{-\beta d^*}$ with $\beta = \ln(1/\gamma)$. A larger $\gamma$ reduces variation across temporal distances and may reduce the effective rank of the resulting state-goal matrix.

\section{Contrastive RL: discounted reachability under a logarithmic parametrization and the behavior policy}

At the optimum of the binary noise-contrastive estimation objective, the CRL critic $f(s,a,g)=\ip{\varphi(s,a)}{\psi(g)}$ satisfies (Eysenbach et al.\ 2022, Lem.~4.1) $\exp f^* = p^{\pi}*\gamma(s*+{=}g \mid s,a)/p(g)$, which is a discounted state-occupancy ratio. Under deterministic dynamics, suppose a trajectory generated by the behavior policy first visits $g$ at step $T^{\pi}(s,g)$ and does not revisit it. Then $p^{\pi}*\gamma(g|s) = (1-\gamma)\gamma^{T^{\pi}(s,g)}$, and
\begin{equation}
f^*(s,g) ;=; -\ln!\tfrac{1}{\gamma}; T^{\pi}(s,g) ;+; c(g), \qquad c(g) = \ln(1-\gamma) - \ln p(g).
\label{eq:crl}
\end{equation}
Comparing \eqref{eq:factorization} and \eqref{eq:crl}, \textbf{both methods use a rank-$d$ bilinear parametrization for quantities derived from $\gamma^{,T(s,g)}$, but they use different transformations and correspond to different policies}. CRL represents $\log K^{\pi} + c(g)$, using a logarithmic parametrization of the behavior-policy hitting time with a goal-dependent additive term. GCIQL represents $K^*$ through an affine transformation of the optimal hitting time, without a goal-dependent additive term. Myers et al.\ (2024) recover a quasimetric from occupancy critics through self-subtraction, $d*{\mathrm{SD}}(s,g) := f(g,g)-f(s,g)$, which removes $c(g)$. Applied to GCIQL value estimates, where $V^*(g,g)=0$, self-subtraction gives $-V^* = D_\gamma$. Thus, for bilinear GCIQL, the construction of Myers et al.\ and the negated optimal value coincide, and Prop.~\ref{prop:qm} shows that the result is a quasimetric when the value function is optimal. FB representations provide a more general factorization: $F^\top B$ factorizes the successor measure, whose goal-state marginal ratio yields the quantity in \eqref{eq:crl}. GCIQL's $K^*$ is the corresponding optimal-policy case for the sparse goal-reaching reward.

\section{VIP: the symmetric positive-semidefinite case yields a Euclidean embedding}

VIP specifies $V(s;g) = -\lVert \phi(s)-\phi(g)\rVert_2$ and trains $\phi$ so that the embedding distance represents a KL-regularized optimal value. The connection to bilinear GCIQL follows from linear algebra. Suppose the learned factorization is approximately symmetric and positive semidefinite on the support of the offline dataset, so that $K(s,g)=\ip{u(s)}{u(g)}$ for a shared embedding $u$. Equivalently, the symmetric component of $K$ is dominant and positive semidefinite, in which case such an embedding exists by Mercer's theorem. Using $K(s,s) = 1 \iff V(s,s)=0$, polarization gives
\begin{equation}
\big\lVert u(s) - u(g) \big\rVert_2^2 ;=; K(s,s) + K(g,g) - 2K(s,g) ;=; 2\big(1-\gamma^{,d^*(s,g)}\big) ;=; 2(1-\gamma),\big({-}V^*(s,g)\big).
\label{eq:polarization}
\end{equation}
\textbf{A symmetric positive-semidefinite bilinear GCIQL solution therefore defines an embedding of the form used by VIP}: squared Euclidean distance in the embedding space equals $2(1-\gamma)$ times the negated value. VIP uses $\lVert\cdot\rVert$, whereas GCIQL implies $\lVert\cdot\rVert^2$. Since $\sqrt{\cdot}$ is concave and satisfies $\sqrt{0}=0$, $\lVert u(s)-u(g)\rVert = \sqrt{2(1-\gamma)D_\gamma}$ is also a metric-preserving transformation, and $u$ defines a metric embedding into $\R^{d+1}$. Schoenberg's theorem gives the converse characterization: $e^{-\beta d}$ is positive semidefinite for all $\beta>0$ if and only if $d$ is conditionally negative definite, also called a metric of negative type or a Hilbert-embeddable metric. Therefore, bilinear GCIQL can represent a symmetric positive-semidefinite solution for every $\gamma$ if and only if the optimal goal-reaching distance is of negative type. The antisymmetric component of the learned $K$ measures deviation from this condition and is evaluated by diagnostic P4 in the experiment specification. Finally, if embedding norms are approximately constant, $\lVert u \rVert \approx c$, then \eqref{eq:polarization} gives $V \approx c^2\big(\cos\angle(u(s),u(g)) - 1\big)/(1-\gamma)$ up to constants. Thus, the statement that ``GCIQL learns cosine similarity'' applies only when the learned factorization is approximately symmetric and positive semidefinite and the embedding norms are concentrated, as evaluated by diagnostic P5.

\section{QRL: an alternative parametrization of the same target quantity}
QRL learns $d_\theta \approx d^*$ directly, enforces the quasimetric axioms through the IQE architecture, and replaces temporal-difference learning with a constrained maximization objective. Prop.~\ref{prop:qm} shows that bilinear GCIQL targets a bounded transformation of the same optimal goal-reaching distance through a different parametrization. The bilinear parametrization is unconstrained, and a low-rank approximation does not necessarily preserve the composition inequality in Cor.~\ref{cor:mult}. Therefore, $\hat K$ is not guaranteed to define a valid discounted goal-reachability function outside the support of the training distribution. This gives a diagnostic that does not require ground-truth goal distances:

\begin{proposition}[Triangle-inequality diagnostic]\label{prop:cert}
Let $\hat d(s,g) := \log_\gamma \hat K(s,g)$, where $\hat K = 1 + (1-\gamma)\hat V$ is clipped to $(0,1]$. If $\hat V$ is the optimal value function of any goal-reaching MDP, then $\hat d$ satisfies the triangle inequality. Therefore, any observed triangle-inequality violation implies $\hat V \notin {V^* : \text{goal-reaching MDPs}}$, without requiring ground-truth distances.
\end{proposition}

\section{Summary and testable predictions}

All four methods can be related through the discounted reachability quantity $K = e^{-\beta d}$ induced by a quasimetric $d$.
\begin{center}
\resizebox{\textwidth}{!}{%
\begin{tabular}{l l l l}
\hline
Method & Target quantity & Parametrization & Policy / constraint \\
\hline
Bilinear GCIQL & $K^* = \gamma^{d^*}$ & affine & optimal over dataset-supported actions; unconstrained rank $d$ \\
CRL & $\log K^{\pi} + c(g)$ & logarithmic & behavior policy; unconstrained rank $d$ \\
VIP & $-\phi_\gamma(d^*)$ (KL-reg.) & $-\lVert\phi(s)-\phi(g)\rVert$ & regularized optimal policy; symmetric by construction \\
QRL & $d^*$ & direct distance parametrization & optimal policy; quasimetric constraints enforced architecturally \\
\hline
\end{tabular}}
\end{center}
These identities correspond directly to the diagnostics in the companion experiment specification. Equation~\eqref{eq:factorization} motivates comparing the singular-value spectrum and effective rank of $\hat K$ with the performance transition across representation dimensions (P7). Equation~\eqref{eq:crl} motivates measuring rank correlation between GCIQL's $-\hat V$ and CRL's $f(g,g)-f(s,g)$ (P6). Equation~\eqref{eq:polarization} predicts that, when the learned solution is approximately symmetric and positive semidefinite, squared distances $\lVert u(s)-u(g)\rVert^2$ in the embedding space should induce rankings similar to those produced by $-\hat V$ and the VIP distance (P4--P6). Norm concentration determines whether cosine similarity provides an adequate approximation (P5). Proposition~\ref{prop:cert} motivates measuring triangle-inequality violations of $\hat d$ as a diagnostic of whether the learned critic is consistent with an optimal goal-conditioned value function (P3).

\paragraph{Caveats.} Equations~\eqref{eq:vstar}--\eqref{eq:crl} are exact under deterministic dynamics when trajectories visit the goal only once. Under stochastic dynamics, $\gamma^{d^*}$ is replaced by $\E[\gamma^{T}]$, and hitting-time expressions are replaced by discounted occupancy measures. OGBench cube is deterministic when the full state is observed, so these identities apply on the support of the offline dataset. An expectile parameter of $\tau=0.9$ approximates but does not attain the maximum over dataset-supported actions. The condition $V(g,g)=0$ holds at the fixed point but is not explicitly enforced; consequently, the diagonal estimate $\hat V(s,s)$ provides a simple convergence diagnostic. Finally, $d^*$ denotes reachability under policies supported by the offline dataset, not unrestricted reachability in the underlying MDP.





\printbibliography[title={References}]

\end{document}
